% part: intuitionistic-logic
% chapter: soundness-completeness
% section: canonical-model

\documentclass[../../../include/open-logic-section]{subfiles}

\begin{document}

\olfileid{int}{sc}{mod}

\olsection{النموذج القانوني}

ستكون العوالم في نموذجنا متتاليات منتهية~$\sigma$ من الأعداد الطبيعية؛
أي~$\sigma \in \Nat^*$. ولاحظ أن~$\Nat^*$ معرّفة استقرائيًا كما يأتي:
\begin{enumerate}
\item $\emptyseq \in \Nat^*$.
\item إذا كان~$\sigma \in \Nat^*$ و$n \in \Nat$، كان
  $\sigma.n \in \Nat^*$ (حيث~$\sigma.n$ هي
  $\sigma \concat \tuple{n}$، و$\sigma \concat \sigma'$ هو وصل
  $\sigma$ و$\sigma'$).
\item لا شيء آخر ينتمي إلى~$\Nat^*$.
\end{enumerate}
ومن ثم يمكننا استعمال~$\Nat^*$ في وضع تعريفات استقرائية.

لتكن~$\tuple{!B_1, !C_1}$، $\tuple{!B_2, !C_2}$، \dots، تعدادًا
لجميع أزواج ال!!{formula}s. وبالنسبة إلى مجموعة
ال!!{formula}s~$\Delta$، نعرّف~$\Delta(\sigma)$ استقرائيًا كما يأتي:
\begin{enumerate}
\item $\Delta(\emptyseq) = \Delta$.
\item $\Delta(\sigma.n) = {}$
  \[
  \begin{cases}
    (\Delta(\sigma) \cup \{!B_n\})^* &
    \text{إذا كان $\Delta(\sigma) \cup \{!B_n\} \Proves/ !C_n$} \\
    \Delta(\sigma) & \text{خلاف ذلك}
  \end{cases}
  \]
\end{enumerate}
ونعني بـ~$(\Delta(\sigma) \cup \{!B_n\})^*$ هنا مجموعة
ال!!{formula}s الأولية التي يضمن وجودها تطبيق
\olref[lin]{lem:lindenbaum} على المجموعة
$\Delta(\sigma) \cup \{!B_n\}$ وال!!{formula}~$!C_n$. ولاحظ أنه بحسب
هذا التعريف، إذا كان~$\Delta(\sigma) \cup \{!B_n\} \Proves/ !C_n$،
كان~$\Delta(\sigma.n) \Proves !B_n$ و
$\Delta(\sigma.n) \Proves/ !C_n$. ولاحظ أيضًا أن
$\Delta(\sigma) \subseteq \Delta(\sigma.n)$ لأي~$n$. وإذا كانت
$\Delta$ أولية، كانت~$\Delta(\sigma)$ أولية لكل~$\sigma$.

\begin{defn}\ollabel{defn:canonical-model}
  افترض أن~$\Delta$ أولية. عندئذ يُعرّف \emph{النموذج القانوني}
  $\mModel{M(\Delta)}$ لـ~$\Delta$ كما يأتي:
  \begin{enumerate}
  \item $W = \Nat^*$، وهي مجموعة المتتاليات المنتهية من الأعداد الطبيعية.
  \item $R$ هو الترتيب الجزئي الذي فيه~$R\sigma\sigma'$ إذا وفقط إذا
    كانت~$\sigma$ قطعةً ابتدائية من~$\sigma'$ (أي
    $\sigma'=\sigma\concat\sigma''$ لمتتالية ما~$\sigma''$).
  \item $V(p) = \Setabs{\sigma}{p \in \Delta(\sigma)}$.
  \end{enumerate}
\end{defn}

يسهل التحقق من أن~$R$ ترتيب جزئي بالفعل. ويتحقق أيضًا شرط رتابة~$V$.
فمن~$\Delta(\sigma)\subseteq\Delta(\sigma.n)$ نحصل بالاستقراء على
$\Delta(\sigma)\subseteq\Delta(\sigma')$ كلما كان
$R\sigma\sigma'$.

\end{document}
